\section{Limitations}
\label{sec:limitations}
\vspace{-1ex}\stitle{Synthetic data realism.}
Our experiments on perturbed real tables validate the observed degradation on generated tables. However, \name is designed as a controlled diagnostic generator, not as a replacement for manually curated benchmarks built from naturally occurring documents.
Although the generator is conditioned on domain descriptions, unit inventories, and optional examples, the resulting tables may not capture all distributional, visual, and semantic regularities of real spreadsheets or reports, not even with \name applied in a semi-automatic manner (\S\ref{app:real-constraint}).
This is a deliberate trade-off: from-scratch generation gives us control over table count, reasoning topology, perturbations, and units, but it does not guarantee full realism. 

\stitle{Restricted numerical structure.} This study focuses on numerical QA, where each latent table contains a distinguished numerical value attribute. Although \name{} can also generate non-numerical questions, we restrict our analysis to numerical questions because parallel scenarios, which require evidence retrieval from multiple tables followed by an aggregation operation, are especially common with numerical values. Future work could extend \name{} to non-numerical settings, including questions that require set operations, categorical comparisons, or logical aggregation over textual evidence.

% This study focuses on numerical QA where each latent table contains a distinguished numerical value attribute. Although \name{} could be used to generate non-numerical questions, we restrict this study on numerical questions, as parallel scenarios, involving evidence retrieval followed by an aggregation operation, are more frequent with numbers. Future work [continue]
% This design simplifies pivoting, unit conversion, provenance tracking, and executable answer generation, but it is narrower than many real tables, which may contain multiple interacting numerical measures. \mlc{è vero che ogni rappresentazione latente contiene un valore float, ma è solo una decisione sulla tipologia di domande: siamo partiti con lo sviluppo delle domande parallele, e queste domande sono per definizione numeriche (retrieval di evidence + aggregazione). Il sistema può essere usato anche per domande di lookup (come fatto nelle domande sequenziali, dove le tabelle iniziali non contengono valori float), o esteso per altre tipologie di domande, mantenendo sempre l'executable answer generation, o anche informazioni di provenance in alcuni casi}

\stitle{Question verbalization.}
Gold answers are computed by executing latent SQL programs, and non-relational views are generated through answer-preserving transformations.
However, the NL question is produced by an LLM and is therefore not guaranteed correct by construction. We mitigate this issue through verifier-based repair and manual inspection.
\autoref{tab:correctness} shows high reliability overall, but also that sequential questions are more error-prone, mainly because verbalization may drop latent selection constraints. Thus, \name provides executable guarantees for answers and evidence preservation, while question fidelity remains only an empirically validated component.

\stitle{Generation failures and filtering.}
Some generations fail schema-format checks. We discard invalid generations, which improves benchmark quality but increases the cost of data generation. %\S\ref{sec:data_statistics} reports the dropout rates.
Future work should focus on reducing the number of discarded generations to maximize QA generation throughput.
% Some generations fail schema-format checks and some perturbations are only applicable when the latent table satisfies structural constraints. We discard invalid generations, which improves benchmark quality but may bias the final distribution toward easier-to-render schemas. Future work should study how filtering affects the diversity of generated domains, layouts, and reasoning patterns. \mlc{inserisco in appendice le informazioni sul dropout rate così da quantificare gli errori di generazione. Inserisco qui un link all'appendice}

\stitle{Perturbation coverage.}
We consider a broad set of structural, cell-level, layout, and numerical perturbations. Nevertheless, real non-relational tables contain additional phenomena that we do not fully model, such as visual alignment cues and chart-table mixtures. Although \name{} could be extended to cover those perturbations, the current benchmark does not stress the full range of document-table understanding.  We leave this as future work.

% We consider a broad set of structural, cell-level, layout, and numerical perturbations. Nevertheless, real non-relational tables contain additional phenomena that we do not fully model, such as visual alignment cues and chart-table mixtures. The current benchmark therefore stresses important non-relational phenomena, but not the full range of document-table understanding. \mlc{anche se nelle limitazioni, inserirei qualcosa sul fatto che il set di perturbazioni è estendibile senza compromettere l'accuratezza del sistema, seguita da una frase su sviluppi futuri.}

% \stitle{Metric limitations.} We use normalized exact match against scalar gold answers. This metric is appropriate for automatically verifiable numerical QA, but it may under-credit semantically equivalent answers when models produce rounded values, alternative units, or explanatory responses that contain the correct value in a non-canonical form. \mlc{forse toglierei quest'ultima: i prompt utilizzati indicano il formato di output, il numero di valori decimali, e le domande indicano sempre l'unità di misura che si desidera nella risposta. L'inferenza, essendo in CoT, produce sempre reasoning traces (explanatory responses) seguite dalla risposta finale... quindi empiricamente, nel nostro setting, non ci sono problemi di metrica, che invece potrebbero presentarsi se le risposte fossero non numeriche}